\chapter{P\&L and the Accounting of a Position}\label{ch:m1:pnl-and-positions}

At 16:05 the trader tells her desk head she made forty thousand dollars. At
17:30 the risk report says thirty-one. The next morning the finance
department books twenty-seven. Nobody is lying and nobody has made a mistake:
the three numbers use the same six trades and three different answers to the
questions ``at what price is the \omterm{def:m1:pnl-and-positions:position}{position} worth what?'' and ``what did the
day cost?''. A trading firm lives and pays its people on this number. This
chapter defines it, shows which part of it is arithmetic and which part is
convention, and builds the component that computes it.

\section{Positions and exposure}

\begin{definition}[Position]\label{def:m1:pnl-and-positions:position}
A \emph{position}\index{position} in an instrument is the signed quantity
held: positive (\emph{long}) if owned, negative (\emph{short}) if owed. It
changes only through fills, corporate actions and transfers; it is the sum of
the signed quantities of all of them since inception.
\end{definition}

\begin{definition}[Notional, gross and net exposure]\label{def:m1:pnl-and-positions:exposure}
The \emph{notional}\index{notional} of a \omterm{def:m1:pnl-and-positions:position}{position} $q_i$ marked at price $p_i$
is $|q_i|\,p_i$ (times the contract multiplier, for a derivative). For a
book, the \emph{gross exposure}\index{gross exposure} is $\sum_i |q_i| p_i$
and the \emph{net exposure}\index{net exposure} is $\sum_i q_i p_i$.
\end{definition}

\begin{example}[A long--short book]\label{ex:m1:pnl-and-positions:book}
Long \$300 million and short \$200 million on \$100 million of equity: gross
\$500 million, or $5\times$ the equity, the measure of how much can go wrong;
net \$100 million, or $1\times$, the measure of how much market direction is
being taken. Limits are set on both, and on each \omterm{def:m1:pnl-and-positions:position}{position} separately.
\end{example}

\section{Marking to market}

\begin{definition}[Mark-to-market]\label{def:m1:pnl-and-positions:mtm}
To \emph{mark to market}\index{mark-to-market} a \omterm{def:m1:pnl-and-positions:position}{position} is to value it at a
current market price, the \emph{mark}, instead of at what was paid for it.
The profit and loss of a book over a period is the change in its marked value
plus the net cash it produced.
\end{definition}

\begin{proposition}[The P\&L identity]\label{prop:m1:pnl-and-positions:identity}
Start flat. After fills $k = 1,\dots,n$ with sides $\epsilon_k = \pm1$ (buy
$+1$), quantities $n_k$, prices $p_k$ and fees $f_k$, the \omterm{def:m1:pnl-and-positions:position}{position} is $q =
\sum_k \epsilon_k n_k$, the cash is $c = -\sum_k \epsilon_k n_k p_k$, and the
profit and loss at mark $m$ is
\[
\pnl(m) \;=\; c + q\,m - \sum_k f_k .
\]
It depends on the fills only through $(q, c, \sum f_k)$: on no accounting
convention, and not on the order of the fills.
\end{proposition}
\begin{proof}
The book holds cash $c - \sum f_k$ and $q$ units worth $m$ each; it started
with nothing.
\end{proof}

This is the exact part. Everything is an integer if prices are in ticks or
ledger units (\cref{bld:m1:what-a-trading-firm-does:ledger}), and a P\&L
system that cannot reproduce it to the unit has lost a trade. The
\emph{split} of that number is where conventions enter.

\begin{definition}[Average cost, realised and unrealised P\&L]\label{def:m1:pnl-and-positions:realised}
Under the \emph{average cost}\index{average cost} convention an open \omterm{def:m1:pnl-and-positions:position}{position}
carries the quantity-weighted average price $\bar p$ of the fills that built
it. A fill that reduces the \omterm{def:m1:pnl-and-positions:position}{position} by $n$ units at price $p$
\emph{realises} $\pm n(p - \bar p)$ (plus for a long, minus for a short) and
leaves $\bar p$ unchanged; a fill that increases it updates $\bar p$; a fill
that flips it closes the old \omterm{def:m1:pnl-and-positions:position}{position} entirely and opens the remainder at
$p$. The \emph{realised P\&L}\index{realised P\&L} is the sum of these
amounts; the \emph{unrealised P\&L}\index{unrealised P\&L} is $q(m - \bar
p)$.
\end{definition}

\begin{example}[Six fills]\label{ex:m1:pnl-and-positions:six}
The trader of the opening paragraph buys 20\,000 at 49.30 and 20\,000 at
49.55 ($\bar p = 49.425$), sells 10\,000 at 50.10 (realising $10\,000 \times
0.675 = 6\,750$), buys 10\,000 at 49.70 ($\bar p = 49.494$), sells 25\,000 at
50.28 (realising $19\,656$), and buys 15\,000 at 50.20 ($\bar p = 49.847$).
She ends long 30\,000 with \$26\,406 realised. Under \emph{first in, first
out} the same sales are matched to the oldest purchases and realise
\$28\,750. Marked at 50.00 the unrealised parts are \$4\,594 and \$2\,250:
both conventions total \$31\,000, as \cref{prop:m1:pnl-and-positions:identity}
says they must. The split matters for tax and for performance statistics
(``win rate''); it never changes what the firm is worth.
\end{example}

\begin{omfigure}
\begin{tikzpicture}
\begin{axis}[width=11cm,height=5.2cm,
  xlabel={},ylabel={\$ thousand},
  xtick={0,1,2,3,4,5,6,7},xticklabels={09:30,09:35,10:10,11:20,13:05,14:30,15:40,16:00},
  tick label style={font=\footnotesize},label style={font=\small},
  xmin=0,xmax=7,ymin=-2,ymax=42,grid=major,grid style={gray!25},
  legend columns=3,legend style={font=\footnotesize,at={(0.5,-0.22)},anchor=north,draw=none,fill=none,/tikz/every even column/.append style={column sep=8pt}},
  table/col sep=comma]
\addplot[omMeth,very thick,const plot] table[x=k,y=realised_k]{figdata/markets-1/07-pnl-and-positions/day.csv};
\addplot[omProp,very thick,mark=*,mark size=1.2pt] table[x=k,y=unrealised_k]{figdata/markets-1/07-pnl-and-positions/day.csv};
\addplot[omDef,very thick,mark=*,mark size=1.4pt] table[x=k,y=total_k]{figdata/markets-1/07-pnl-and-positions/day.csv};
\legend{realised (average cost),unrealised,total at the mid}
\end{axis}
\end{tikzpicture}

\omcaption{The day of \cref{ex:m1:pnl-and-positions:six}, observed at each
fill and at the close, marked at the mid. \omterm{def:m1:pnl-and-positions:realised}{Realised P\&L} moves only when a
sale reduces the \omterm{def:m1:pnl-and-positions:position}{position}; the last point falls because the closing mid is
below the afternoon's prices while she is still long 30\,000 shares. Data:
the chapter's fill list.}\label{fig:m1:pnl-and-positions:day}
\end{omfigure}

\begin{method}[Choosing a mark]\label{met:m1:pnl-and-positions:mark}
\begin{enumerate}
\item \textbf{The last trade} is the most recent information and the worst
  mark: it is at the bid or at the ask, at someone else's size, possibly
  minutes old.
\item \textbf{The mid} of a tight, two-sided quote is the usual intraday mark
  for liquid instruments.
\item \textbf{The side against you} --- bid for longs, ask for shorts ---
  values the \omterm{def:m1:pnl-and-positions:position}{position} at what closing it would fetch for a small size; for a
  large \omterm{def:m1:pnl-and-positions:position}{position} subtract an estimate of impact as well.
\item \textbf{The official close} (closing auction price or settlement price)
  is what finance, clearing houses and fund administrators use: it is
  public, unique and auditable.
\end{enumerate}
Decide once, in writing, which mark serves which report. A trader allowed to
choose the mark can choose the P\&L.
\end{method}

\begin{example}[Forty, thirty-one, twenty-seven]\label{ex:m1:pnl-and-positions:three}
The trader's screen marks at the last print, 50.30: cash $-\$1\,469\,000$
plus $30\,000 \times 50.30$ gives \$40\,000. The risk system marks at the
closing mid, 50.00: \$31\,000. The \$9\,000 between them is $30\,000$ shares
$\times$ 30 cents of mark. Finance uses the official close, also 50.00, and
subtracts \$3\,400 of fees and \$600 of financing: \$27\,000
(\cref{fig:m1:pnl-and-positions:three}).
\end{example}

\begin{omfigure}
\begin{tikzpicture}
\begin{axis}[width=11cm,height=3.6cm,xbar,bar width=8pt,
  xlabel={P\&L of the same six fills (\$ thousand)},
  ytick=data,yticklabels from table={figdata/markets-1/07-pnl-and-positions/three_numbers.csv}{label},
  yticklabel style={font=\small},tick label style={font=\small},label style={font=\small},
  y dir=reverse,xmin=0,xmax=45,enlarge y limits=0.25,grid=major,grid style={gray!25},
  nodes near coords,nodes near coords style={font=\footnotesize},
  table/col sep=comma]
\addplot[fill=omDef!60,draw=omDef] table[x=usd_k,y=k]{figdata/markets-1/07-pnl-and-positions/three_numbers.csv};
\end{axis}
\end{tikzpicture}

\omcaption{One day, three numbers. Data: \cref{ex:m1:pnl-and-positions:three},
computed by the chapter's script.}\label{fig:m1:pnl-and-positions:three}
\end{omfigure}

\section{What a position costs and earns while it is held}

\begin{definition}[Carry]\label{def:m1:pnl-and-positions:carry}
The \emph{carry}\index{carry} of a \omterm{def:m1:pnl-and-positions:position}{position} is the P\&L it produces if prices
do not move: income received (dividends, coupons, interest on short proceeds)
less costs paid (financing of longs, borrow fees on shorts, dividends owed on
shorts).
\end{definition}

A complete daily P\&L therefore has five lines besides price moves:
commissions and exchange fees on the day's fills; financing accrued
(\cref{bld:m1:financing:calc}); borrow fees; dividends and coupons, booked on
the ex-date (\cref{ch:m1:shares-corporate-actions-indices}); and, for
\omterm{def:m1:pnl-and-positions:position}{positions} in a foreign currency, the move of the exchange rate.

\begin{proposition}[Currency translation]\label{prop:m1:pnl-and-positions:fx}
A \omterm{def:m1:pnl-and-positions:position}{position} of $q$ units of a foreign asset moves from price $p_0$ to $p_1$
in its own currency while the exchange rate (base currency per unit of
foreign currency) moves from $x_0$ to $x_1$. Its P\&L in base currency is
\[
q\,(p_1x_1 - p_0x_0) \;=\; \underbrace{q\,(p_1-p_0)\,x_0}_{\text{price}}
 \;+\; \underbrace{q\,p_0\,(x_1-x_0)}_{\text{currency}}
 \;+\; \underbrace{q\,(p_1-p_0)(x_1-x_0)}_{\text{cross}} .
\]
\end{proposition}
\begin{proof}
Expand $p_1x_1 = (p_0 + \Delta p)(x_0 + \Delta x)$.
\end{proof}

\begin{example}[A good stock in a bad currency]\label{ex:m1:pnl-and-positions:fx}
A dollar-based fund holds 10\,000 shares of a euro stock that rises from 40
to 41 while the euro falls from 1.10 to 1.08 dollars. Price: $+\$11\,000$.
Currency: $-\$8\,000$. Cross: $-\$200$. Total $+\$2\,800$. A fund that wants
the first number and not the second sells \euro 400\,000 forward; the hedge
covers the opening value $qp_0$, never the cross term, and must be resized as
the stock moves.
\end{example}

\section{P\&L explain: a first look}

\begin{definition}[P\&L attribution]\label{def:m1:pnl-and-positions:attribution}
A \emph{P\&L attribution}\index{P\&L attribution}, or \emph{explain}, is a
decomposition of a period's P\&L into named causes whose sum equals the
total, with any remainder reported as \emph{unexplained}.
\end{definition}

\begin{method}[The daily explain of a cash book]\label{met:m1:pnl-and-positions:explain}
With $q_0$ the \omterm{def:m1:pnl-and-positions:position}{position} at the previous close $m_0$, and $m_1$ today's close:
\begin{enumerate}
\item \textbf{Carried \omterm{def:m1:pnl-and-positions:position}{position}:} $q_0(m_1 - m_0)$.
\item \textbf{New trades:} $\sum_k \epsilon_k n_k (m_1 - p_k)$ --- each fill
  measured from its own price to the close.
\item \textbf{Fees}, \textbf{financing}, \textbf{dividends},
  \textbf{currency}: each its own line.
\item \textbf{Unexplained:} total less the sum of the above. For a cash book
  it should be zero to the cent; anything else is a missing trade, a wrong
  mark or an unprocessed corporate action.
\end{enumerate}
\end{method}

The identity behind the method is
\cref{prop:m1:pnl-and-positions:identity} applied between two closes: the
first two lines sum to $c + q_1 m_1 - q_0 m_0$. Line 2 is where a trading
desk's skill shows: a \omterm{def:m1:what-a-trading-firm-does:market-maker}{market maker}'s new-trade P\&L is its spread capture and
its \omterm{def:m1:what-a-trading-firm-does:adverse-selection}{adverse selection} (\cref{prop:m1:what-a-trading-firm-does:capture}); its
\omterm{def:m1:pnl-and-positions:position}{carried-position} P\&L should be small and is pure risk. For derivatives the
first line is itself split by risk factor, which is the subject of One Quant
Book 6.

\begin{omfigure}
\begin{tikzpicture}
\begin{axis}[width=11cm,height=4.4cm,xbar,bar width=8pt,
  xlabel={\$ thousand},
  ytick=data,yticklabels from table={figdata/markets-1/07-pnl-and-positions/explain.csv}{label},
  yticklabel style={font=\small},tick label style={font=\small},label style={font=\small},
  y dir=reverse,xmin=-4,xmax=14,enlarge y limits=0.15,grid=major,grid style={gray!25},
  nodes near coords,nodes near coords style={font=\footnotesize},
  table/col sep=comma]
\addplot[fill=omDef!60,draw=omDef] table[x=usd_k,y=k]{figdata/markets-1/07-pnl-and-positions/explain.csv};
\end{axis}
\end{tikzpicture}

\omcaption{The next day's explain: 30\,000 shares carried from 50.00 to 50.40,
10\,000 sold at 50.35 and 5\,000 bought at 50.10, \$900 of fees and \$650 of
financing. Most of the day's \$11\,450 was made by yesterday's \omterm{def:m1:pnl-and-positions:position}{position}, not
by today's decisions. Data: computed by the chapter's script.}
\end{omfigure}

\begin{omfigure}
\begin{tikzpicture}[font=\small,
  box/.style={draw=omDef,rounded corners=2pt,minimum width=2.6cm,minimum height=0.9cm,align=center,fill=omDef!4},
  src/.style={draw=omExo,rounded corners=2pt,minimum width=2.6cm,minimum height=0.9cm,align=center,fill=omExo!6},
  flow/.style={-{Stealth[length=2mm]},thick,omDef}]
\node[src] (fills) at (0,1.2) {Fills\\(drop copy)};
\node[src] (marks) at (0,-1.2) {Marks\\(policy per report)};
\node[box] (keeper) at (4.6,0) {Position and\\P\&L keeper};
\node[box] (explain) at (9.2,1.2) {Explain};
\node[box] (ledger) at (9.2,-1.2) {Ledger};
\node[src] (recon) at (13.4,0) {Reconciliation\\with the broker};
\draw[flow] (fills) -- (keeper);
\draw[flow] (marks) -- (keeper);
\draw[flow] (keeper) -- (explain);
\draw[flow] (keeper) -- (ledger);
\draw[flow] (ledger) -- (recon);
\draw[flow] (explain) -- (recon);
\end{tikzpicture}

\omcaption{Where the keeper sits in the miniature firm. Fills arrive from the
venue's or broker's drop copy, never from the strategy's own belief about
what it traded; \omterm{def:m1:pnl-and-positions:position}{positions} and cash are reconciled every morning against the
clearing broker's statement.}
\end{omfigure}

\section{Tutorial: three numbers, one day}\label{tut:m1:pnl-and-positions:three}

\begin{tutorial}
\textbf{Goal.} Replay the six fills, reproduce \$40\,000, \$31\,000 and
\$27\,000, and check that the accounting convention does not matter.
\textbf{End state:} \cref{fig:m1:pnl-and-positions:day,fig:m1:pnl-and-positions:three}.
\begin{tutsteps}
\item \textbf{The keeper's core}, as it is written for the systems side of
the firm. Quantity, cash and fees are integers; only the reported split is
floating point.
\omcode{code/firm/pnl/cpp/firm_pnl.hpp}{21}{45}{Average-cost position keeping in C++: add, reduce, flip; and the exact total.}\label{lst:m1:pnl-and-positions:cpp}
\item \textbf{Replay and mark three ways.}
\omcode{code/markets-1/07-pnl-and-positions/python/pnl_day.py}{28}{41}{The same position object, three marks and two cost lines.}\label{lst:m1:pnl-and-positions:three}
\item \textbf{Run the tests}: \texttt{make test-code CH=markets-1/07-pnl-and-positions}.
They assert the three numbers, and that first-in-first-out gives the same
total with a different split (\$28\,750 realised against \$26\,406).
\item \textbf{Stress the identity.} The build's tests replay five thousand
random fills and require the split to agree with the exact total.
\end{tutsteps}
\textbf{What to change next.} Mark the long at the bid, 49.99: what is the
fourth number? Then add a dividend of 20 cents going ex during the day and
decide in which line of the explain it belongs.
\end{tutorial}

\section{Build: the position and P\&L keeper}\label{bld:m1:pnl-and-positions:keeper}

\begin{build}
\bfield{Purpose} The component every strategy, risk check and report of the
miniature firm asks: what do we hold, and what is it worth? It is the first
piece built in all three languages, because it runs both inside the
low-latency process (One Quant Book 13) and in the research and reporting
stack.
\bfield{Interface} \texttt{Position.on\_fill(side, quantity, price, fee)},
\texttt{unrealised(mark)}, \texttt{total(mark)}; \texttt{Book.on\_fill(symbol,
\dots)}, \texttt{total(marks)}, \texttt{gross\_exposure(marks)},
\texttt{net\_exposure(marks)}. Prices, cash and fees are integers in ledger
units.
\bfield{Rules} \texttt{total} is computed from quantity, cash and fees only,
and is exact. The realised and unrealised split follows
\cref{def:m1:pnl-and-positions:realised}. Non-positive quantities or prices
and negative fees are rejected. No allocation and no exception on the fill
path in the C++ and Rust versions other than for rejected input.
\bfield{Acceptance tests} \texttt{code/firm/pnl/tests/}, \texttt{cpp/} and
\texttt{rust/}: round trip; \omterm{def:m1:pnl-and-positions:realised}{average cost} and partial close; flip; five
thousand random fills with split equal to total; rejections. The three
implementations must agree on the same fill file.
\bfield{Stretch} Add first-in-first-out lots behind the same interface and a
switch per account; post each day's explain lines to the ledger of
\cref{ch:m1:what-a-trading-firm-does}.
\end{build}

\begin{omsources}
\item L.~Harris, \emph{Trading and Exchanges}, Oxford University Press, 2003,
chapter 21 (performance evaluation).
\item R.~Grinold and R.~Kahn, \emph{Active Portfolio Management}, 2nd ed.,
McGraw-Hill, 2000, chapter 17 (performance analysis).
\item International Accounting Standards Board, \emph{IFRS 13 Fair Value
Measurement} --- the accounting standard behind the choice of marks.
\end{omsources}

\section{Exercises}

\begin{exercise}[$\star$]\label{exo:m1:pnl-and-positions:1}
A book is long 40\,000 shares at \$25, long 10\,000 at \$120 and short
30\,000 at \$50, on equity of \$2 million. Give its gross and \omterm{def:m1:pnl-and-positions:exposure}{net exposure} in
dollars and as multiples of equity.
\end{exercise}

\begin{exercise}[$\star$]\label{exo:m1:pnl-and-positions:2}
Starting flat: buy 500 at 20.00, buy 300 at 20.40, sell 600 at 20.50. Give
the \omterm{def:m1:pnl-and-positions:realised}{average cost} after each fill, the \omterm{def:m1:pnl-and-positions:realised}{realised P\&L}, and the \omterm{def:m1:pnl-and-positions:realised}{unrealised P\&L}
at a mark of 20.30.
\end{exercise}

\begin{exercise}[$\star$]\label{exo:m1:pnl-and-positions:3}
For the fills of the previous exercise compute the cash, and check
\cref{prop:m1:pnl-and-positions:identity} at the mark 20.30.
\end{exercise}

\begin{exercise}[$\star\star$]\label{exo:m1:pnl-and-positions:4}
Starting flat: sell 1\,000 at 80.00, sell 1\,000 at 81.00, buy 3\,000 at
79.50. Track the \omterm{def:m1:pnl-and-positions:position}{position}, the \omterm{def:m1:pnl-and-positions:realised}{average cost} and the \omterm{def:m1:pnl-and-positions:realised}{realised P\&L} through the
flip, and give the total at a mark of 79.80.
\end{exercise}

\begin{exercise}[$\star\star$]\label{exo:m1:pnl-and-positions:5}
A sterling-based fund holds 50\,000 shares of a US stock. Over a month the
stock goes from \$200 to \$190 and the dollar from 0.80 to 0.84 pounds.
Compute the price, currency and cross terms and the total in pounds.
\end{exercise}

\begin{exercise}[$\star\star$]\label{exo:m1:pnl-and-positions:6}
Yesterday's close was 30.00 with a \omterm{def:m1:pnl-and-positions:position}{position} of $-20\,000$ shares. Today the
desk buys 8\,000 at 29.60 and sells 5\,000 at 29.90; the close is 29.70; fees
are \$260, the borrow fee accrued is \$120 and the short proceeds earned
\$65 of interest. Produce the explain and the closing \omterm{def:m1:pnl-and-positions:position}{position}.
\end{exercise}

\begin{exercise}[$\star\star\star$]\label{exo:m1:pnl-and-positions:7}
\emph{Coding.} Implement first-in-first-out for \omterm{def:m1:pnl-and-positions:position}{positions} that may be short
as well as long (the chapter's \texttt{fifo} handles a long-only day). Test
it on the fills of exercise~\ref{exo:m1:pnl-and-positions:4} and report the
\omterm{def:m1:pnl-and-positions:realised}{realised P\&L}; compare with \omterm{def:m1:pnl-and-positions:realised}{average cost}.
\end{exercise}

\begin{exercise}[$\star\star\star$]\label{exo:m1:pnl-and-positions:8}
\emph{Find the flaw.} A desk reports a win rate of 78\% ``on closed trades''
and a positive \omterm{def:m1:pnl-and-positions:realised}{realised P\&L} for eleven consecutive months; its total P\&L
over the period is negative. Explain how both can be true, which convention
makes it easy, and which single report would have revealed it in the first
month.
\end{exercise}

\section{Problem: Three Numbers, One Day}

\begin{problem}[{Weekend problem --- reconciling a desk's P\&L}]\label{pb:m1:pnl-and-positions:1}
You are the desk's new quant. On your first morning three P\&L figures for
yesterday are on the head of desk's table and he wants one. The desk trades
one stock. It started the day long 12\,000 shares; the previous close was
84.00. Yesterday's fills, in order: sell 7\,000 at 84.60; sell 9\,000 at
84.90; buy 6\,000 at 84.20; buy 10\,000 at 83.70. The last trade of the day
printed at 83.55, the closing quote was 83.70 bid, 83.80 ask, and the closing
auction price was 83.76. Fees were 0.12 cent a share; financing and borrow
cost \$310.

\medskip\noindent\textbf{Part I --- \omterm{def:m1:pnl-and-positions:position}{Position} and cash.}
\begin{enumerate}
\item Give the \omterm{def:m1:pnl-and-positions:position}{position} after each fill and at the close.
\item Give the cash produced by the day's fills.
\item Compute the fees.
\item At which moments was the desk short, and how many shares?
\end{enumerate}

\medskip\noindent\textbf{Part II --- Three marks.}
\begin{enumerate}[resume]
\item Give the day's gross P\&L marked at the last trade.
\item At the closing mid.
\item At the closing auction price.
\item At the price that would be obtained by closing the \omterm{def:m1:pnl-and-positions:position}{position} against
the quote.
\item Which mark would you use for the trader's bonus, and why?
\end{enumerate}

\medskip\noindent\textbf{Part III --- Explain} (use the auction price).
\begin{enumerate}[resume]
\item Compute the \omterm{def:m1:pnl-and-positions:position}{carried-position} line.
\item Compute the new-trades line, fill by fill.
\item Add fees and financing and give the net total.
\item Check that the unexplained is zero.
\item The trader says ``I made money trading and lost it on the overnight
\omterm{def:m1:pnl-and-positions:position}{position}''. Is that what the explain says?
\end{enumerate}

\medskip\noindent\textbf{Part IV --- Realised and unrealised.}
\begin{enumerate}[resume]
\item The opening 12\,000 shares have an \omterm{def:m1:pnl-and-positions:realised}{average cost} of 82.50. Track the
\omterm{def:m1:pnl-and-positions:realised}{average cost} and the \omterm{def:m1:pnl-and-positions:realised}{realised P\&L} through the four fills.
\item Give the \omterm{def:m1:pnl-and-positions:realised}{unrealised P\&L} at the auction price.
\item Realised plus unrealised is not the day's P\&L. What is it, and what
must be subtracted to recover the day's figure?
\item The clearing broker's statement shows a closing \omterm{def:m1:pnl-and-positions:position}{position} of 13\,000
shares. List three possible causes in order of likelihood and what you would
check first.
\item State the \emph{named result}: the single net P\&L for the day that
you give the head of desk, and its three largest components.
\item In one sentence, state the policy that would have prevented three
numbers from reaching his table.
\end{enumerate}
\end{problem}

\section{Interview questions}

\begin{interviewq}[$\star$ \iqroles{trader, developer, researcher}]\label{iq:m1:pnl-and-positions:1}
What is the difference between realised and \omterm{def:m1:pnl-and-positions:realised}{unrealised P\&L}? Does the choice
between first-in-first-out and \omterm{def:m1:pnl-and-positions:realised}{average cost} change how much money you made?
\end{interviewq}

\begin{interviewq}[$\star$ \iqroles{developer}]\label{iq:m1:pnl-and-positions:2}
Why should a P\&L system not store prices and cash as floating-point numbers?
What do you use instead?
\end{interviewq}

\begin{interviewq}[$\star\star$ \iqroles{trader, researcher}]\label{iq:m1:pnl-and-positions:3}
You are long 10\,000 shares, the last trade is 50.30 and the market is 49.99
bid, 50.01 ask. What is your \omterm{def:m1:pnl-and-positions:position}{position} worth? What if the \omterm{def:m1:pnl-and-positions:position}{position} were two
million shares?
\end{interviewq}

\begin{interviewq}[$\star\star$ \iqroles{developer, trader}]\label{iq:m1:pnl-and-positions:4}
Your strategy believes it is flat; the broker's end-of-day file says it is
long 400 shares. Walk through how you find out which is right, and what the
system should have done the moment they diverged.
\end{interviewq}

\begin{interviewq}[$\star\star$ \iqroles{researcher, bank}]\label{iq:m1:pnl-and-positions:5}
What is a P\&L explain, and what does a persistent unexplained amount tell
you?
\end{interviewq}

\begin{interviewq}[$\star\star\star$ \iqroles{developer}]\label{iq:m1:pnl-and-positions:6}
Design the data model of a \omterm{def:m1:pnl-and-positions:position}{position} keeper that must handle fills arriving
out of order, busted (cancelled) trades and price corrections. What is
stored, and what is derived?
\end{interviewq}
